% Punkte und Strecken im Koordinatensystem · Darstellen und Lage lesen · Prüfungs-Fokus Abitur GK – bank/punkte-und-strecken-im-koordinatensystem/e1.jsonl (zusammenbau v0.6, Prüfungs-Fokus)
\einheitenkopf[e1]{Einheit 1 · Punkte darstellen und Lage lesen}
\zweigzeile{Abi GK ×6}
\setcounter{aufgabe}{0}

% Anlauf (ohne Prüfkennung)
\begin{aufgabe}{Darstellen und Lage lesen – Anlauf}
\begin{teile}
\teil $P(4 | -1 | 0)$ – was sagt die Null? \\ \kreuz{$P$ liegt auf der $x_3$-Achse} \\ \kreuz{$P$ liegt in der $x_1x_2$-Ebene} \\ \kreuz{$P$ liegt in der $x_2x_3$-Ebene}
\teil $A(5 | 3 | 1)$ und $B(2 | 1 | 4)$ eintragen, $C$ ablesen.

\begin{ksys3} \rpunkt{4}{5}{2}{C} \end{ksys3}
C( \leerfeld | \leerfeld | \leerfeld )
\teil Zwei Masten stehen senkrecht in $F_1(-2 | -2 | 0)$ ($4$ m hoch) und $F_2(5 | 3 | 0)$ ($3$ m hoch); ein Seil verbindet die Spitzen (1 LE = 1 m). Zeichne Masten und Seil.

\begin{ksys3} \end{ksys3}
\end{teile}
\end{aufgabe}

\begin{aufgabe}{Darstellen und Lage lesen – Prüfungsaufgaben}
\begin{teile}
\teil Ein Holzwürfel mit Kante $4$ hat eine Ecke im Ursprung und liegt im ersten Oktanten (alle Koordinaten $\ge 0$). Zeichne das Viereck $P(4 | 0 | 1)$, $Q(1 | 4 | 0)$, $R(0 | 4 | 2)$, $S(3 | 0 | 4)$ ein. \hfill (A19)

\begin{ksys3} \rquader{0,0,0}{4}{4}{4} \end{ksys3}
\teil Ein Glaswürfel mit Kante $3$ hat eine Ecke im Ursprung und liegt im ersten Oktanten (alle Koordinaten $\ge 0$). Zeichne das Viereck $P(3 | 1 | 0)$, $Q(3 | 3 | 2)$, $R(1 | 3 | 3)$, $S(0 | 0 | 2)$ ein. \hfill (A19)

\begin{ksys3} \rquader{0,0,0}{3}{3}{3} \end{ksys3}
\teil Ein Würfel mit Kante $4$ hat eine Ecke im Ursprung und liegt im ersten Oktanten (alle Koordinaten $\ge 0$). Zeichne das Viereck $P(0 | 3 | 0)$, $Q(4 | 4 | 1)$, $R(4 | 1 | 4)$, $S(1 | 0 | 4)$ ein. \hfill (A19)

\begin{ksys3} \rquader{0,0,0}{4}{4}{4} \end{ksys3}
\teil Ein Quader hat eine Ecke im Ursprung $O$ und die Kanten $5$, $6$ und $3$ entlang der $x_1$-, $x_2$- und $x_3$-Achse. $P(5 | 2 | 3)$ liegt auf einer Kante. Zeichne das Dreieck mit den Ecken $O$, $P$ und $R(0 | 6 | 3)$ ein. \hfill (A26)

\begin{ksys3} \rquader{0,0,0}{5}{6}{3} \end{ksys3}
\teil Ein Quader hat eine Ecke im Ursprung $O$ und die Kanten $6$, $3$ und $4$ entlang der $x_1$-, $x_2$- und $x_3$-Achse. $P(6 | 0 | 1)$ liegt auf einer Kante. Zeichne das Dreieck mit den Ecken $O$, $P$ und $T(0 | 3 | 4)$ ein. \hfill (A26)

\begin{ksys3} \rquader{0,0,0}{6}{3}{4} \end{ksys3}
\teil Ein Quader hat eine Ecke im Ursprung $O$ und die Kanten $3$, $6$ und $3$ entlang der $x_1$-, $x_2$- und $x_3$-Achse. $P(3 | 2 | 3)$ liegt auf einer Kante. Zeichne das Dreieck mit den Ecken $O$, $P$ und $T(3 | 6 | 0)$ ein. \hfill (A26)

\begin{ksys3} \rquader{0,0,0}{3}{6}{3} \end{ksys3}
\end{teile}
\end{aufgabe}

\begin{aufgabe}{Darstellen und Lage lesen – Prüfungsaufgaben – weiter}
\begin{teile}
\teil Begründe: $P(-3 | 4 | 3)$ und $Q(5 | -1 | 2)$ liegen auf derselben Seite der $x_1x_2$-Ebene. \hfill (A25)
\teil Begründe: Das Dreieck $A(6 | 0 | 1)$, $B(6 | 3 | 5)$, $C(6 | -3 | 5)$ liegt parallel zur $x_2x_3$-Ebene und symmetrisch zur $x_1x_3$-Ebene. \hfill (A26)
\teil Begründe: $R(1 | -2 | -3)$ und $T(4 | 5 | 2)$ liegen auf verschiedenen Seiten der $x_1x_2$-Ebene. \hfill (A25)
\teil Der Grundriss eines Pavillons ist ein Achteck in der $x_1x_2$-Ebene, symmetrisch zur $x_1$- und zur $x_2$-Achse. Gegeben sind $A(6 | 0)$, $B(5 | 3)$, $C(0 | 4)$. Zeichne das vollständige Achteck. \hfill (A25)

\begin{ksys}[xmin=-7,xmax=7,ymin=-7,ymax=7,xlabel=$x_1$,ylabel=$x_2$] \punkt{6}{0}{A} \punkt{5}{3}{B} \punkt{0}{4}{C} \end{ksys}
\teil Der Beet ist ein Achteck in der $x_1x_2$-Ebene, symmetrisch zur $x_1$- und zur $x_2$-Achse. Gegeben sind $A(4 | 0)$, $B(3 | 2)$, $C(0 | 3)$. Zeichne das vollständige Achteck. \hfill (A25)

\begin{ksys}[xmin=-7,xmax=7,ymin=-7,ymax=7,xlabel=$x_1$,ylabel=$x_2$] \punkt{4}{0}{A} \punkt{3}{2}{B} \punkt{0}{3}{C} \end{ksys}
\teil Der Bodenplatte ist ein Achteck in der $x_1x_2$-Ebene, symmetrisch zur $x_1$- und zur $x_2$-Achse. Gegeben sind $A(6 | 0)$, $B(2 | 5)$, $C(0 | 6)$. Zeichne das vollständige Achteck. \hfill (A25)

\begin{ksys}[xmin=-7,xmax=7,ymin=-7,ymax=7,xlabel=$x_1$,ylabel=$x_2$] \punkt{6}{0}{A} \punkt{2}{5}{B} \punkt{0}{6}{C} \end{ksys}
\end{teile}
\end{aufgabe}

\begin{aufgabe}{Darstellen und Lage lesen – Prüfungsaufgaben – weiter}
\begin{teile}
\teil Eine Pyramide hat den quadratischen Boden $ABCD$ mit der Seite $3$ cm; die Spitze $S$ liegt $4$ cm senkrecht über $A$. Im Gitter sind der Boden und die Spitze des Seitendreiecks $ABS$ markiert. Zeichne die drei übrigen Seitendreiecke ein. \hfill (A24)

\begin{ksys}[xmin=-5,xmax=9,ymin=-5,ymax=9] \punkt{0}{0}{A} \punkt{3}{0}{B} \punkt{3}{3}{C} \punkt{0}{3}{D} \punkt{0}{-4}{S} \end{ksys}
\teil Eine Pyramide hat den quadratischen Boden $ABCD$ mit der Seite $4$ cm; die Spitze $S$ liegt $3$ cm senkrecht über $A$. Im Gitter sind der Boden und die Spitze des Seitendreiecks $ABS$ markiert. Zeichne die drei übrigen Seitendreiecke ein. \hfill (A24)

\begin{ksys}[xmin=-4,xmax=10,ymin=-4,ymax=10] \punkt{0}{0}{A} \punkt{4}{0}{B} \punkt{4}{4}{C} \punkt{0}{4}{D} \punkt{0}{-3}{S} \end{ksys}
\teil Eine Pyramide hat den quadratischen Boden $ABCD$ mit der Seite $2$ cm; die Spitze $S$ liegt $1{,}5$ cm senkrecht über $A$. Im Gitter sind der Boden und die Spitze des Seitendreiecks $ABS$ markiert. Zeichne die drei übrigen Seitendreiecke ein. \hfill (A24)

\begin{ksys}[xmin=-2.5,xmax=5.5,ymin=-2.5,ymax=5.5] \punkt{0}{0}{A} \punkt{2}{0}{B} \punkt{2}{2}{C} \punkt{0}{2}{D} \punkt{0}{-1.5}{S} \end{ksys}
\teil Die Pyramide $ABCS$ hat die Grundfläche $A(4 | 1 | 2)$, $B(1 | 5 | 2)$, $C(-1 | -1 | 2)$ und die Spitze $S(4 | -1 | 6)$. Eingezeichnet sind die Projektionen $A'$ und $B'$ in die $x_1x_2$-Ebene. Ergänze $C'$ und $S'$ und entscheide: Liegt der Höhenfußpunkt innerhalb der Grundfläche? \hfill (A18)

\begin{ksys3} \rpunkt*{4}{1}{0}{A'} \rpunkt*{1}{5}{0}{B'} \end{ksys3}
\teil Die Pyramide $ABCS$ hat die Grundfläche $A(5 | -1 | 1)$, $B(1 | 5 | 1)$, $C(-1 | -2 | 1)$ und die Spitze $S(1 | 1 | 6)$. Eingezeichnet sind die Projektionen $A'$ und $B'$ in die $x_1x_2$-Ebene. Ergänze $C'$ und $S'$ und entscheide: Liegt der Höhenfußpunkt innerhalb der Grundfläche? \hfill (A18)

\begin{ksys3} \rpunkt*{5}{-1}{0}{A'} \rpunkt*{1}{5}{0}{B'} \end{ksys3}
\end{teile}
\end{aufgabe}

\medskip
\uebersichtskasten{Drei Koordinaten, drei Wege: P(3 | −2 | 5) heißt vom Ursprung drei Einheiten in x-Richtung, zwei entgegen der y-Richtung, fünf nach oben; im Schrägbild ist die nach vorn weisende Achse schräg und verkürzt gezeichnet – Längen misst man nicht im Bild, man zählt Koordinatenwege. \\ Lage ablesen: Koordinate null → der Punkt liegt in der zugehörigen Koordinatenebene; gleiche Koordinate bei allen Punkten → die Figur liegt parallel zu einer Koordinatenebene; unterscheiden sich zwei Punkte nur im Vorzeichen einer Koordinate, liegen sie symmetrisch zu dieser Koordinatenebene; gleiches Vorzeichen → dieselbe Seite. \\ A(3 | −2 | 5) und B(3 | 2 | 5) liegen symmetrisch zur xz-Ebene; C(3 | −2 | 0) liegt in der xy-Ebene. \\ Projektion: parallel zu einer Achse in die Koordinatenebene projizieren heißt, die zugehörige Koordinate null setzen; die übrigen bleiben. \\ P(3 | −2 | 5) → P′(3 | −2 | 0) in der xy-Ebene.}
